\section{Veri ve Kapsam}

Bu rapor Cumhuriyet Halk Partisi'nin 2015 ve 2026 yılları arasında Kemal 
Kılıçdaroğlu ve Özgür Özel tarafından yapılan konuşmalarını (n = 2.703\footnote{Ham konuşma sayısı = 2.744}) iki farklı nicel metin analizi yaklaşımıyla analiz etmektedir. 
Sözcük düzeyindeki analizde 6 farklı sözlüğün 9 farklı kırılma anında nasıl davrandığı,
bağlam düzeyindeki analizdeyse 11 yıl boyunca \textit{millet}, \textit{kadın} ve \textit{erdoğan} kelimelerinin
anlam kümeleri incelenmiştir.

Raporda analiz edilen konuşmalar \href{https://chp.org.tr}{CHP'nin internet sitesinden} otomatik olarak çekilmiş ve
veri toplama esnasında araştırmanın amacı\footnote{London School of Economics and Political Science (LSE) Metodoloji Departmanı, Uygulamalı Sosyal Veri Bilimi yüksek lisans programı bitirme tezi.} user-agent kullanılarak tanımlanmıştır. Ayrıca her toplanan metin arasında
sunucuya yük bindirmemek için makul (5 saniye) bir gecikme süresi uygulanmıştır. 

Korpus, etkinlik (873), basın (800), miting (445), meclis (347), televizyon (214), ve kurultay (24)
olmak üzere 6 farklı mekandaki konuşmalardan bir araya getirilmiştir\footnote{Birden fazla başlıkta paylaşılan konuşmalar ayıklandı}.
Miting konuşmalarının coğrafi dağılımı [TABLO X]'te paylaşılmıştır.
13 Ocak 2015 ve 8 Kasım 2023 arasında Kemal Kılıçdaroğlu'nun 1.569; 9 Kasım 2023 ve 31 Mart 2026 arasında Özgür Özel'in 1.134 konuşması tespit edilmiştir. 
Görev süresince Özgür Özel daha fazla basın ve miting; Kemal Kılıçdaroğlu ise daha fazla meclis ve etkinlik konuşması gerçekleştirmiştir.
Mekan yoğunluğunun hem sözcük hem bağlamsal bulguları etkileyip etkilemediğine yönelik dayanıklılık testleri yapılmış ve mekanın bulguları etkilemediği görülmüştür.
Bulguların tespit edilemeyen bir kısmı Özgür Özel ve Kemal Kılıçdaroğlu arasındaki Türkçe kullanımından kaynaklanıyor olabilir. Örneğin, Kılıçdaroğlu'na kıyasla Özel daha
çeşitli bir kelime dağarcığı tercih etmekte ve cümle başına yaklaşık 2 kelime daha fazla kullanmaktadır. Buna ek olarak rapor boyunca liderlerin retorik kullanımlarına yönelik bulgulara da yer verilmiştir.


Bu araştırmanın kapsamında yalnızca parti liderlerinin konuşmalarına odaklanılmıştır. Milletvekillerinin açıklamaları için
partinin internet sitesinde seçilmiş konuşmalar yayınlanmakta olup tam kapsamlı bir korpus için Türkiye Büyük Millet Meclisi arşivine başvurulabilir.
Ekrem İmamoğlu'nun konuşmalarının analiz edilmesi planlansa da yeterli veri kaynağına erişilememiştir. Liderlerin sosyal medya 
paylaşımları internet sitesinden erişilebilir olsa da yeterli hacimde bir veri kaynağı oluşturmadığından analiz dışında bırakılmıştır.

Dokümanların nasıl temsil edildiği, korpus temizlik aşamaları, kullanılan yöntemler ve algoritmalar yöntem bölümünde açıklanmıştır. 
Frekans ve kırılma analizi için oluşturulan sözlükler yöntem bölümünde mevcuttur.
Analizi yeniden üretmek veya kodu ve dayanıklılık testlerini incelemek isteyen araştırmacılar e.avci@lse.ac.uk adresinden iletişime geçebilir.




